% for glossary entry
% @entry{bird,
%     name={bird},                          % Nome completo e scorciatoia per richiamare il termine nel testo
%     description = {feathered animal},     % Spiegazione per la pagina del glossrio 
%     see={[see also]{duck,goose}}          % Se sono presenti termini correlati all'acronimo
% }

% if this bib file does not work, try using \input{file.tex}
% where all the \newabbreviation commands have been inserted
% containing all the definitions

% Possibili comandi:
% \gls{parola} permette di stampare il termine (la prima volta esteso e con l'abbreviazione tra parentesi, le successive volte solo l'abbreviazione)
% \Gls{parola} permette di stampare il termine con la prima lettera maiuscola (la prima volta esteso e con l'abbreviazione tra parentesi)
% \GLS{parola} permette di stampare il termine tutto in maiuscolo (la prima volta esteso e con l'abbreviazione tra parentesi)
% \glspl{parola} permette di stampare il termine al plurale
% \glsxtrshort{parola} permette di stampare solo l'abbreviazione (anche se è la prima volta che si usa)


% if you want to use also description for the abbreviations/acronyms, you should use bib2gls and define all the entries in a bib file, which is incompatible with Overleaf
\newacronym{DBMS}{DBMS}{Database Management System}